\documentclass[11pt,a4paper]{article}

% -- Packages ----------------------------------------------------------------
\usepackage[utf8]{inputenc}
\usepackage[T1]{fontenc}
\usepackage{lmodern}
\usepackage[margin=2.5cm]{geometry}
\usepackage{amsmath,amssymb,amsthm}
\usepackage{graphicx}
\usepackage{booktabs}
\usepackage{array}
\usepackage{multirow}
\usepackage{xcolor}
\usepackage{listings}
\usepackage{hyperref}
\usepackage{microtype}
\usepackage{parskip}
\usepackage{caption}
\usepackage{subcaption}
\usepackage{algorithm}
\usepackage{algpseudocode}
\usepackage{enumitem}
\usepackage{url}
\usepackage{cite}

% -- Hyperref setup -----------------------------------------------------------
\hypersetup{
    colorlinks=true,
    linkcolor=blue!70!black,
    citecolor=green!50!black,
    urlcolor=cyan!70!black,
    pdfauthor={Robin Dey},
    pdftitle={Collective Habit as Infrastructure: A Computational Model of
              Morphic Resonance via Multi-Agent Reinforcement Learning},
    pdfsubject={Morphic Resonance, Multi-Agent Systems, Reinforcement Learning,
                Collective Memory, Stigmergy},
    pdfkeywords={morphic resonance, Q-learning, multi-agent systems,
                 collective memory, displaced control, stigmergy,
                 knowledge transfer, cold start}
}

% -- Listings / code blocks ---------------------------------------------------
\definecolor{codebg}{HTML}{F5F5F5}
\definecolor{codeframe}{HTML}{CCCCCC}
\definecolor{keywordcolor}{HTML}{0033BB}
\definecolor{commentcolor}{HTML}{008800}
\definecolor{stringcolor}{HTML}{BB0000}

\lstset{
    backgroundcolor=\color{codebg},
    frame=single,
    rulecolor=\color{codeframe},
    basicstyle=\ttfamily\footnotesize,
    keywordstyle=\color{keywordcolor}\bfseries,
    commentstyle=\color{commentcolor}\itshape,
    stringstyle=\color{stringcolor},
    breakatwhitespace=false,
    breaklines=true,
    captionpos=b,
    keepspaces=true,
    numbersep=5pt,
    showspaces=false,
    showstringspaces=false,
    showtabs=false,
    tabsize=2,
    xleftmargin=0.5em,
    xrightmargin=0.5em,
    aboveskip=0.5em,
    belowskip=0.5em
}

\lstdefinelanguage{Python}{
    keywords={class, def, return, if, else, elif, for, while, import, from,
              as, in, not, and, or, None, True, False, self, lambda, with,
              try, except, raise, finally, yield, assert, pass, break,
              continue, is, del, global, nonlocal},
    morecomment=[l]{\#},
    morestring=[b]{"""},
    morestring=[b]{"},
    morestring=[b]{'},
    sensitive=true
}

% -- Math helpers -------------------------------------------------------------
\DeclareMathOperator*{\argmax}{arg\,max}
\DeclareMathOperator*{\argmin}{arg\,min}
\newtheorem{definition}{Definition}

% -- Title block --------------------------------------------------------------
\title{\textbf{Collective Habit as Infrastructure:\\
       A Computational Model of Morphic Resonance\\
       via Multi-Agent Reinforcement Learning}\\[0.5em]
       \large Technical Report --- BUTTERS Project}

\author{Robin Dey \\
        \textit{BUTTERS / GraphPalace Project} \\
        \href{https://github.com/web3guru888/GraphPalace}{\texttt{github.com/web3guru888/GraphPalace}}}

\date{April 2026}

% -----------------------------------------------------------------------------
\begin{document}
\maketitle

% -- Abstract -----------------------------------------------------------------
\begin{abstract}
We present a computational model of Rupert Sheldrake's morphic resonance
hypothesis, implemented as a multi-agent reinforcement learning simulation of
the McDougall rat maze experiment (1920--1954).  The simulation places
populations of tabular Q-learning agents in a procedurally generated maze
across 60 generations under four experimental conditions: no inheritance
(control), genetic inheritance only, morphic field only, and both mechanisms
combined.  The morphic field is modeled as a performance-weighted shared
Q-table that accumulates knowledge from all past agents and emits it as a
scaled prior for new agents, with no genetic or physical channel of
transmission.  A fifth condition --- the \emph{displaced control} --- drops
fresh agents with zero ancestry into a mature morphic field at generation~60.
Results show: (1)~the control condition remains flat at 400~steps across all
generations; (2)~genetic-only converges to 56~steps by generation~5 (86\%
improvement); (3)~morphic-only improves to 245~steps (39\%) with no genetic
inheritance; (4)~displaced control agents immediately score 297~steps versus
the 400-step baseline --- a 26\% advantage on their first generation with
zero ancestry.  We argue that the simulation does not demonstrate the
existence of morphic resonance in nature, but establishes three things: the
hypothesis is formally coherent, the predicted data pattern is distinctive and
detectable, and the displaced control constitutes the most powerful
experimental design for testing it.  We further show that the morphic field
architecture --- performance-weighted absorption, scaled emission, and
blending with local learning --- constitutes a practical design pattern for
multi-agent knowledge transfer, cold-start mitigation, and stigmergic
coordination in AI systems.

\vspace{0.5em}
\noindent\textbf{Keywords:} morphic resonance, Q-learning, multi-agent
systems, collective memory, displaced control, stigmergy, knowledge transfer,
cold start problem
\end{abstract}

\tableofcontents
\newpage

% =============================================================================
\section{Introduction}
\label{sec:intro}
% =============================================================================

Rupert Sheldrake's hypothesis of morphic resonance~\cite{sheldrake1981,
sheldrake1988} proposes that natural systems inherit a collective memory from
all previous systems of the same kind, transmitted not through known physical
mechanisms but through \emph{morphic fields} --- non-material regions of
influence that shape form, development, and behavior.  The hypothesis has been
a lightning rod for controversy since its publication: \emph{Nature} editor
John Maddox called Sheldrake's first book ``the best candidate for
burning''~\cite{maddox1981}, while supporters argue that orthodox biology
lacks an account of how identical genomes produce radically different cell
types, or why the so-called ``laws'' of nature should be eternal rather than
habitual~\cite{sheldrake2012}.

The hypothesis remains untested in a decisive way, in part because it is
difficult to formalize.  What \emph{exactly} is a morphic field?  How does
similarity-based resonance operate?  What data pattern would we expect to see
if it were real, and how would we distinguish that pattern from genetic
selection, environmental drift, or statistical artifact?

This paper takes a computational approach.  We do not attempt to prove or
disprove morphic resonance.  Instead, we ask: \emph{if we build the
hypothesis into a multi-agent simulation, what happens?}  Specifically:

\begin{enumerate}[leftmargin=2em]
    \item Can a ``morphic field,'' formalized as a shared, performance-weighted
          Q-table with no genetic channel, transmit useful learning across
          agent generations?
    \item What is the distinctive data pattern that would result --- the
          \emph{signature} an experimentalist should look for?
    \item Does the displaced control test (fresh, unrelated agents dropped into
          a mature field) constitute a viable experimental design?
    \item Does the resulting architecture have practical value for multi-agent
          AI systems, independent of Sheldrake's metaphysics?
\end{enumerate}

Our simulation is built on the McDougall rat maze experiment
(1920--1954)~\cite{mcdougall1938}, chosen because it is the most quantitative,
generational, and independently replicated of the six major evidence categories
for morphic resonance (Section~\ref{sec:background:mcdougall}).  We model rats
as tabular Q-learning agents, generations as population cohorts, genetic
inheritance as Q-table crossover with mutation, and the morphic field as an
accumulated, performance-weighted Q-table that emits a scaled prior to
unrelated agents.

\paragraph{Contributions.}
\begin{enumerate}[leftmargin=2em]
    \item A formal computational model of morphic resonance using
          multi-agent reinforcement learning (Section~\ref{sec:architecture}).
    \item Empirical results across five experimental conditions demonstrating
          the displaced control signature
          (Section~\ref{sec:results}).
    \item An epistemological analysis of what the simulation does and does
          not prove (Section~\ref{sec:epistemology}).
    \item Identification of the morphic field as a practical design pattern
          for multi-agent knowledge transfer, cold-start mitigation, and
          stigmergic coordination (Section~\ref{sec:applications}).
\end{enumerate}

\paragraph{Paper organization.}
Section~\ref{sec:background} reviews Sheldrake's hypothesis, the McDougall
experiment, and related work in multi-agent learning.
Section~\ref{sec:architecture} describes the simulation architecture.
Section~\ref{sec:results} presents experimental results.
Section~\ref{sec:epistemology} analyzes what the results mean.
Section~\ref{sec:applications} connects the morphic field to practical AI
design patterns.  Section~\ref{sec:discussion} discusses limitations.
Section~\ref{sec:conclusion} concludes.


% =============================================================================
\section{Background and Related Work}
\label{sec:background}
% =============================================================================

% -- 2.1 Morphic Resonance ---------------------------------------------------
\subsection{Morphic Resonance}
\label{sec:background:morphic}

A \emph{morphic field} (from Greek \emph{morphe}, form) is, in Sheldrake's
framework, a field of information that shapes the form, development, and
behavior of a system~\cite{sheldrake1981}.  Morphic fields are localized
within and around \emph{morphic units} --- any unit of organization, from
subatomic particles to organisms to social groups --- and are organized in a
nested hierarchy (a \emph{holarchy}, per Koestler~\cite{koestler1967}).

\emph{Morphic resonance} is the process by which information is transmitted
from past morphic units to present ones on the basis of \emph{similarity}: the
more similar a present system is to past systems, the stronger the resonance.
The mechanism is explicitly non-energetic --- no photons, no spacetime
curvature, no known carrier --- and non-local in both space and time.
Sheldrake frames this as a radical departure from the assumption that the
``laws'' of nature are fixed and eternal, proposing instead that they are
\emph{habits} that accumulate through repetition~\cite{sheldrake1988}.

\paragraph{Key terminology.}
\textbf{Morphic unit}: any system with its own morphic field.
\textbf{Morphogenetic field}: the subset governing embryonic development
(the original concept from Waddington and others, which Sheldrake
generalizes).  \textbf{Behavioral field}: governs learned or instinctive
behavior.  \textbf{Chreode}: a canalized pathway of development
(Waddington's term, adopted by Sheldrake for habitual trajectories in morphic
fields).

\paragraph{How morphic fields differ from known physical fields.}
(A)~No known energy is transmitted.  (B)~No inverse-square attenuation with
distance.  (C)~Influence is non-local in time: past systems influence present
ones without propagation delay.  (D)~Specificity is by similarity, not
proximity.  These properties make the hypothesis difficult to test but also
generate distinctive predictions.

% -- 2.2 The McDougall Experiment --------------------------------------------
\subsection{The McDougall Rat Maze Experiment}
\label{sec:background:mcdougall}

William McDougall conducted maze-learning experiments with Wistar rats at
Harvard from 1920 to 1954~\cite{mcdougall1938}.  Rats chose between a lighted
gangway (electric shock) and an unlighted gangway (safe).  The first generation
averaged approximately 165~errors; by the 30th generation, this had fallen to
approximately 20~errors.  McDougall interpreted the improvement as Lamarckian
inheritance --- the inheritance of acquired characteristics.

\paragraph{The Agar--Crew replication.}
Agar, Crew, and Mitchell at Edinburgh ran a parallel study from 1934 to 1954,
maintaining both a trained line and, crucially, a \emph{control line} never
selected for maze ability~\cite{agar1954}.  The trained line improved,
consistent with McDougall.  But the control line \emph{also improved at roughly
the same rate}.  The original researchers attributed this to environmental and
methodological artifacts accumulated over 20~years.

\paragraph{Sheldrake's reinterpretation.}
Sheldrake argues that the control-line improvement is the key datum: if rats
never selected for maze ability also improve, something other than genetics is
transmitting the learning~\cite{sheldrake1988}.  This is the \emph{morphic
resonance signature} --- the prediction that our simulation is designed to
test.

\paragraph{Conventional counter-explanations.}
Laboratory condition changes over two decades, genetic drift, unconscious
handling changes, and statistical fluctuation have all been proposed.  No
consensus exists on which explanation is correct.

% -- 2.3 Why This Experiment? ------------------------------------------------
\subsection{Simulatability Ranking}
\label{sec:background:ranking}

Of the six major evidence categories for morphic resonance --- rat maze
learning, crystal formation, blue tit milk-bottle opening, staring detection,
dog telepathy, and telephone telepathy --- we selected the McDougall rat maze
because it is:

\begin{enumerate}[leftmargin=2em]
    \item \textbf{Quantitative and generational}: a clear measurable
          (errors per trial) across successive generations, mapping directly
          to agents and iterations in code.
    \item \textbf{Has the control-line prediction}: the most striking claim
          is not that trained rats' offspring improved, but that
          \emph{unrelated} control lines also improved.
    \item \textbf{Independently replicated}: Agar and Crew ran follow-up
          studies and reportedly found similar trends, though interpretations
          remain disputed.
\end{enumerate}

The remaining evidence categories are weaker for simulation: crystal formation
requires modeling poorly defined physics; blue tit behavior spreading is
adequately explained by network diffusion; telephone telepathy reduces to a
biased coin flip; staring detection is binary guessing susceptible to response
bias~\cite{colwell2000}; and dog telepathy is single-subject observational.

% -- 2.4 Related Work in Multi-Agent Learning --------------------------------
\subsection{Related Work in Multi-Agent Knowledge Transfer}
\label{sec:background:marl}

The idea that agents can benefit from the experience of other agents is well
established in AI, though it is not typically framed in terms of morphic
resonance:

\paragraph{Experience replay and distillation.}
Prioritized experience replay~\cite{schaul2015} allows a single agent to
re-learn from weighted past transitions.  Knowledge
distillation~\cite{hinton2015} compresses the knowledge of a large model (or
ensemble) into a smaller one.  Our morphic field can be viewed as a
population-level experience replay buffer, where contributions are weighted by
performance and ``distilled'' into a shared initialization.

\paragraph{Population-based training.}
Population-based training (PBT)~\cite{jaderberg2017} evolves
hyperparameters and model weights across a population of agents.  Our genetic
condition is a simplified PBT variant; the morphic field adds a
non-genetic channel that PBT does not provide.

\paragraph{Federated learning.}
Federated learning~\cite{mcmahan2017} aggregates locally trained model
updates into a global model without sharing raw data.  The morphic field's
absorb/emit cycle is structurally similar: agents contribute local Q-tables
(weighted by performance) and receive a scaled global prior.

\paragraph{Stigmergy.}
Stigmergy --- indirect coordination through a shared
environment~\cite{grasse1959} --- is the closest analogue.  Ant colony
optimization~\cite{dorigo1996} deposits pheromone on successful paths; our
morphic field deposits weighted Q-values.  The GraphPalace
system~\cite{graphpalace2026} uses five-type pheromone stigmergy for
memory retrieval, operating on the same principle in embedding space rather
than Q-table space.

% -- 2.5 Philosophical Context -----------------------------------------------
\subsection{Philosophical Context}
\label{sec:background:philosophy}

Sheldrake's hypothesis sits within a broader tradition of process philosophy
and holistic biology:

\paragraph{Whitehead.}
Alfred North Whitehead's process philosophy~\cite{whitehead1929} replaces
enduring substances with momentary \emph{actual occasions}, each of which
``prehends'' (grasps) the past.  Sheldrake considers this the most compatible
metaphysical framework for morphic resonance.

\paragraph{Bohm.}
David Bohm's implicate order~\cite{bohm1980} posits a deeper enfolded reality
beneath the manifest (explicate) world.  Morphic fields could be understood as
structures within the implicate order, with resonance as the process by which
enfolded information unfolds into new systems.

\paragraph{Jung.}
Carl Jung's collective unconscious~\cite{jung1959} --- an inherited stratum of
shared archetypes beneath individual consciousness --- is a psychological
parallel.  Morphic resonance generalizes the concept from human psyche to all
of nature.

These frameworks are invoked here not as evidence for morphic resonance, but
as context for the kind of claim being modeled: that collective memory is a
fundamental property of nature, not a special property of brains.


% =============================================================================
\section{Simulation Architecture}
\label{sec:architecture}
% =============================================================================

The simulation (\texttt{morphic\_sim.py}, 493~lines of Python) models the
McDougall experiment as a multi-agent reinforcement learning system with five
experimental conditions.  Figure~\ref{fig:architecture} shows the overall
data flow.

\begin{figure}[ht]
\centering
\begin{verbatim}
  +-------------------+     +-------------------+
  |    Maze (DFS)     |     |   MorphicField    |
  | 7x7 cells, 15x15 |     | shared Q-table    |
  |   grid, seed=42   |     | perf-weighted     |
  +--------+----------+     +--------+----------+
           |                         |
           v                         v emit()
  +-------------------+     +-------------------+
  |   Agent (Q-learn) |<----| Q-table init:     |
  |  lr=0.2, g=0.95   |     | genetic, morphic, |
  |  eps: 0.5 -> 0.05 |     | both, or zeros    |
  +--------+----------+     +--------+----------+
           |                         ^
           | train(100 eps)          | absorb(q, weight)
           | evaluate(greedy)        |
           v                         |
  +-------------------+     +--------+----------+
  | eval_steps (perf) |---->| perf_weight =     |
  | per agent per gen  |     | max_steps / steps |
  +-------------------+     +-------------------+
\end{verbatim}
\caption{Simulation data flow.  Each generation: (1)~agents are initialized
from genetic parents, morphic field, both, or zeros; (2)~agents train for
100~episodes; (3)~agents are evaluated greedily; (4)~Q-tables are absorbed
into the morphic field weighted by performance; (5)~elite parents are selected
for the next generation's genetic channel.}
\label{fig:architecture}
\end{figure}

% -- 3.1 Maze ----------------------------------------------------------------
\subsection{Maze Environment}
\label{sec:architecture:maze}

The maze is generated via recursive backtracker (depth-first search) with
\texttt{seed=42}.  Cell dimensions are $7 \times 7$, yielding a
$15 \times 15$ grid (cells plus walls).  The start position is $(1, 1)$ and
the goal is $(13, 13)$.

\begin{equation}
\label{eq:maze}
\texttt{grid}[r, c] =
\begin{cases}
1 & \text{wall} \\
0 & \text{open passage}
\end{cases}
\end{equation}

The maze is deterministic: all experimental conditions share the same maze
topology, isolating the effect of the knowledge-transfer mechanism.

% -- 3.2 Agent ---------------------------------------------------------------
\subsection{Q-Learning Agent}
\label{sec:architecture:agent}

Each agent is a tabular Q-learning agent operating on the grid:

\begin{itemize}[leftmargin=2em]
    \item \textbf{State space}: $(r, c)$ positions on the $15 \times 15$ grid.
    \item \textbf{Action space}: 4 cardinal directions (up, down, left, right).
    \item \textbf{Q-table}: $15 \times 15 \times 4$ array of float64 values.
    \item \textbf{Hyperparameters}: learning rate $\alpha = 0.2$, discount
          factor $\gamma = 0.95$, $\epsilon$ decays linearly from 0.5 to 0.05
          over training episodes.
    \item \textbf{Rewards}: $+100$ for reaching goal, $-1$ per step, $-5$
          for wall collisions.
    \item \textbf{Training}: 100 episodes per agent, max 400 steps per
          episode.
    \item \textbf{Evaluation}: greedy ($\epsilon = 0$), returns steps to goal
          or 400 (max) if unsolved.
\end{itemize}

The Q-learning update rule is standard:

\begin{equation}
\label{eq:qlearn}
Q(s, a) \leftarrow Q(s, a) + \alpha \left[ r + \gamma \max_{a'} Q(s', a') - Q(s, a) \right]
\end{equation}

% -- 3.3 Morphic Field -------------------------------------------------------
\subsection{Morphic Field}
\label{sec:architecture:field}

The morphic field is the core contribution.  It is modeled as a
\texttt{MorphicField} class holding an accumulated Q-table of the same shape
as agent Q-tables ($15 \times 15 \times 4$).

\begin{definition}[Morphic Field]
A morphic field $\mathcal{F}$ is a tuple $(\mathbf{T}, W, n)$ where
$\mathbf{T} \in \mathbb{R}^{H \times W \times A}$ is the accumulated
Q-table, $W \in \mathbb{R}^+$ is the total weight, and $n \in \mathbb{N}$ is
the contributor count.
\end{definition}

\paragraph{Absorption.}
When an agent completes training and evaluation, its Q-table is absorbed into
the field, weighted by performance:

\begin{equation}
\label{eq:absorb}
\mathbf{T} \leftarrow \mathbf{T} + \mathbf{Q}_i \cdot w_i, \qquad
W \leftarrow W + w_i, \qquad
n \leftarrow n + 1
\end{equation}

where the performance weight is:

\begin{equation}
\label{eq:perfweight}
w_i =
\begin{cases}
\frac{\text{max\_steps}}{\text{eval\_steps}_i} & \text{if agent solved the maze} \\
0.01 & \text{otherwise}
\end{cases}
\end{equation}

This ensures that agents who solve the maze quickly contribute proportionally
more, while agents who fail contribute near-zero weight.  The field does not
fill with noise from failed agents.

\paragraph{Emission.}
When a new agent is initialized from the field, it receives a scaled copy of
the weighted average:

\begin{equation}
\label{eq:emit}
\mathbf{Q}_{\text{init}} = \frac{\mathbf{T}}{W} \cdot s
\end{equation}

where $s = 0.5$ is the emission strength.  This gives new agents a
\emph{nudge} in the direction of past solutions without fully determining
their behavior --- they must still learn.

\paragraph{Key design properties.}
The field has the following properties that map onto Sheldrake's description:
\begin{enumerate}[leftmargin=2em]
    \item \textbf{No genetic channel}: the field operates independently of
          parent--child relationships.  Agents absorb into and emit from the
          field regardless of ancestry.
    \item \textbf{Similarity by task}: all agents solving the same maze
          ``resonate'' with each other through the shared Q-table shape.
    \item \textbf{Cumulative}: the field strengthens with each generation,
          modeling Sheldrake's claim that habits of nature deepen through
          repetition.
    \item \textbf{Performance-weighted}: better solvers contribute more,
          modeling the idea that more successful forms of organization
          leave stronger morphic imprints.
\end{enumerate}

% -- 3.4 Genetic Inheritance --------------------------------------------------
\subsection{Genetic Inheritance}
\label{sec:architecture:genetic}

The genetic channel models standard neo-Darwinian evolution:

\begin{enumerate}[leftmargin=2em]
    \item \textbf{Elite selection}: agents are ranked by evaluation
          performance; the top 30\% (minimum 2) are selected as parents.
    \item \textbf{Crossover}: two random elite parents are selected; a
          uniform binary mask determines which Q-values come from which
          parent.
    \item \textbf{Mutation}: Gaussian noise ($\sigma = 0.05$) is added to
          the child Q-table.
\end{enumerate}

When both genetic and morphic channels are active, the child Q-table is a
50/50 blend:

\begin{equation}
\label{eq:blend}
\mathbf{Q}_{\text{init}} = 0.5 \cdot \mathbf{Q}_{\text{genetic}} + 0.5 \cdot \mathbf{Q}_{\text{morphic}}
\end{equation}

% -- 3.5 Experimental Conditions ---------------------------------------------
\subsection{Experimental Conditions}
\label{sec:architecture:conditions}

Table~\ref{tab:conditions} summarizes the five experimental conditions.

\begin{table}[ht]
\centering
\caption{Experimental conditions.  Each condition enables or disables the
genetic and morphic channels independently.}
\label{tab:conditions}
\begin{tabular}{@{}lccl@{}}
\toprule
\textbf{Condition} & \textbf{Genetic} & \textbf{Morphic} & \textbf{Purpose} \\
\midrule
Control             & \textsf{--} & \textsf{--} & Baseline: no inheritance \\
Genetic Only        & \checkmark  & \textsf{--} & Standard evolution \\
Morphic Field Only  & \textsf{--} & \checkmark  & Collective habit, no genetics \\
Genetic + Morphic   & \checkmark  & \checkmark  & Both mechanisms \\
Displaced Control   & \textsf{--} & \checkmark$^*$ & Fresh agents in mature field \\
\bottomrule
\multicolumn{4}{l}{\footnotesize $^*$Morphic field pre-trained for 60
generations under the morphic-only condition.}
\end{tabular}
\end{table}

The \emph{displaced control} is the critical test.  It corresponds to
Sheldrake's prediction about the Agar--Crew control line: fresh agents with
\textbf{zero} genetic connection to any prior solver are initialized using
only a mature morphic field accumulated over 60 prior generations of
morphic-only training.  If these agents outperform the control baseline on
their first generation, the field alone is transmitting learned information.

% -- 3.6 Pseudocode ----------------------------------------------------------
\subsection{Algorithm}
\label{sec:architecture:algorithm}

Algorithm~\ref{alg:morphic} gives the complete procedure for running one
experimental condition across generations.

\begin{algorithm}[ht]
\caption{Run one experimental condition}
\label{alg:morphic}
\begin{algorithmic}[1]
\Require Maze $M$, condition $\in$ \{control, genetic, morphic, both\},
         generations $G$, population $P$, episodes $E$
\Ensure History $H = [h_1, \ldots, h_G]$ of mean evaluation steps
\State $\mathcal{F} \gets \text{MorphicField}()$ if morphic enabled
\State $\text{parents} \gets \varnothing$
\For{$g \gets 1$ to $G$}
    \For{$i \gets 1$ to $P$}
        \State $\mathbf{Q}_g \gets \varnothing$, $\mathbf{Q}_f \gets \varnothing$
        \If{genetic enabled \textbf{and} parents $\neq \varnothing$}
            \State $\mathbf{Q}_g \gets \text{Breed}(\text{parents})$
        \EndIf
        \If{morphic enabled \textbf{and} $\mathcal{F}.\text{count} > 0$}
            \State $\mathbf{Q}_f \gets \mathcal{F}.\text{emit}(s=0.5)$
        \EndIf
        \State $\mathbf{Q}_{\text{init}} \gets \text{Blend}(\mathbf{Q}_g, \mathbf{Q}_f)$
            \Comment{50/50 if both, else whichever is non-null, else zeros}
        \State $a_i \gets \text{Agent}(M, \mathbf{Q}_{\text{init}})$
        \State $a_i.\text{train}(E)$
    \EndFor
    \State $\text{evals} \gets [a_i.\text{evaluate}() \text{ for } i = 1 \ldots P]$
    \State $h_g \gets \text{mean}(\text{evals})$
    \If{morphic enabled}
        \For{each $(a_i, e_i) \in (\text{agents}, \text{evals})$}
            \State $\mathcal{F}.\text{absorb}(a_i.\mathbf{Q}, w = \text{perf\_weight}(e_i))$
        \EndFor
    \EndIf
    \If{genetic enabled}
        \State $\text{parents} \gets$ top 30\% of agents by eval performance
    \EndIf
\EndFor
\State \Return $H$
\end{algorithmic}
\end{algorithm}


% =============================================================================
\section{Experimental Results}
\label{sec:results}
% =============================================================================

All experiments use seed 42, 60~generations (20~for displaced control),
20~agents per generation, and 100~training episodes per agent.

% -- 4.1 Main Conditions -----------------------------------------------------
\subsection{Main Conditions}
\label{sec:results:main}

Table~\ref{tab:results} presents the headline results.

\begin{table}[ht]
\centering
\caption{Simulation results across conditions (7$\times$7 maze, seed=42,
60~generations, 20~agents/generation, 100~training episodes).  ``Steps''
is the mean evaluation steps in greedy mode (lower is better; 400 = unsolved).}
\label{tab:results}
\begin{tabular}{@{}lrrr@{}}
\toprule
\textbf{Condition} & \textbf{Gen 1 (steps)} & \textbf{Gen 60 (steps)} &
\textbf{Improvement} \\
\midrule
Control             & 400 & 400 &   0\% \\
Genetic Only        & 400 &  56 &  86\% \\
Morphic Field Only  & 400 & 245 &  39\% \\
Genetic + Morphic   & 400 & 108 &  73\% \\
\midrule
Displaced Control$^*$ & 297 & 228 & --- \\
\bottomrule
\multicolumn{4}{l}{\footnotesize $^*$Fresh agents, zero ancestry, mature
morphic field from generation 60.  Runs for 20 generations.}
\end{tabular}
\end{table}

\paragraph{Control.}
Remains flat at 400~steps across all 60 generations.  Individual Q-learning
agents cannot solve this maze within the 100-episode training budget starting
from a zero Q-table.  This establishes a firm baseline.

\paragraph{Genetic Only.}
Converges rapidly to approximately 56~steps by generation~5 and remains stable.
Elite selection with Q-table crossover is a powerful optimization mechanism,
consistent with evolutionary computation results.

\paragraph{Morphic Field Only.}
This is the key morphic resonance result.  Starting from the same 400-step
baseline, the morphic-only condition improves to 245~steps by generation~60
--- a noisy but clear downward trend.  Learning occurs across generations with
\textbf{no genetic inheritance}.  The only channel is the shared morphic field.

The improvement is slower and noisier than the genetic condition.  This is
expected: the morphic field is a blunt instrument compared to elite selection.
It averages over all past agents (performance-weighted), while genetic
inheritance directly copies successful Q-tables with targeted mutation.

\paragraph{Genetic + Morphic.}
Intermediate at 108~steps --- better than morphic-only, worse than
genetic-only.  The 50/50 blending dilutes the genetic signal with the
(noisier) morphic prior in early generations when the field is immature.

% -- 4.2 Displaced Control ---------------------------------------------------
\subsection{The Displaced Control}
\label{sec:results:displaced}

The displaced control is the most important result.  Fresh agents with
\textbf{zero genetic connection} to any prior solver are initialized using
only the mature morphic field (accumulated over 60~generations of morphic-only
training with a separate seed).

\begin{itemize}[leftmargin=2em]
    \item \textbf{Generation 1}: 297~steps (vs.\ 400-step control baseline).
    \item \textbf{Generation 20}: 228~steps.
    \item \textbf{Immediate advantage}: 103~steps (26\%) on the very first
          generation, with zero ancestry.
\end{itemize}

These agents benefit purely because \emph{other agents solved this maze
before}.  There is no genetic inheritance, no training history, no shared
environment --- only the accumulated morphic field.

This directly models Sheldrake's interpretation of the Agar--Crew data:
unrelated rat lines (control groups never selected for maze ability) also
improved over time, suggesting a non-genetic channel of information transfer.


% =============================================================================
\section{Epistemological Analysis}
\label{sec:epistemology}
% =============================================================================

\subsection{What the Simulation Proves}
\label{sec:epistemology:proves}

The simulation demonstrates three things:

\begin{enumerate}[leftmargin=2em]
    \item \textbf{Formal coherence.}  The hypothesis of morphic resonance can
          be formalized as a mathematical model --- specifically, a
          performance-weighted shared Q-table with absorb/emit operations ---
          and this model produces learning across generations with no genetic
          channel.  The hypothesis is not incoherent or self-contradictory.

    \item \textbf{Distinctive data pattern.}  If morphic resonance were real,
          the data pattern would be distinctive: unrelated agents (or
          organisms) solving a task faster purely because other agents solved
          it before, with no physical channel of inheritance.  This pattern is
          detectable and distinguishable from genetic selection (which requires
          ancestry) and individual learning (which cannot exceed the
          single-generation control baseline).

    \item \textbf{Experimental design validation.}  The displaced control test
          is the most powerful experimental design for detecting morphic
          resonance.  If fresh, unrelated agents outperform the control
          baseline on their first generation, something other than genetics
          or individual learning must be transmitting information.
\end{enumerate}

\subsection{What the Simulation Does Not Prove}
\label{sec:epistemology:does_not_prove}

The simulation does \textbf{not} demonstrate:

\begin{enumerate}[leftmargin=2em]
    \item \textbf{That morphic resonance exists in nature.}  In the
          simulation, we \emph{built} the field --- it is a real data
          structure (\texttt{MorphicField} class).  We know exactly how
          information gets from past agents to present ones, because we wrote
          the code.  In Sheldrake's hypothesis, the field's physical
          existence and mechanism are unknown.

    \item \textbf{That information can transmit without a physical channel.}
          The morphic field in the simulation \emph{is} a physical channel
          --- it is a NumPy array in computer memory.  The simulation shows
          what would happen \emph{if} such a channel existed, not that one
          does.

    \item \textbf{That Sheldrake's specific mechanism is real.}  The
          simulation implements \emph{a} formalization of morphic resonance,
          not \emph{the} formalization.  Sheldrake's description is
          qualitative; many formalizations are possible.
\end{enumerate}

\paragraph{The analogy.}
The epistemic status is: ``If gravity worked like an inverse-cube law, orbits
would look like $X$'' (a simulation) versus ``Gravity actually is
inverse-cube'' (an empirical claim).  We have demonstrated the \emph{if}, not
the \emph{is}.

\subsection{The Critical Epistemic Difference}
\label{sec:epistemology:difference}

In the simulation, the morphic field is a \emph{design choice} --- a data
structure we implemented.  In nature, whether any such field exists is an
\emph{empirical question}.  The simulation makes the hypothesis precise enough
to test: run a McDougall-style experiment with rigorous displaced controls; if
the predicted pattern appears (unrelated lines improving), the hypothesis gains
evidence.  If it does not, the hypothesis is weakened.  Either way, the
simulation has served its purpose by sharpening the prediction.


% =============================================================================
\section{Applications to Multi-Agent AI Systems}
\label{sec:applications}
% =============================================================================

Independent of whether morphic resonance exists in nature, the morphic field
architecture --- performance-weighted absorption, scaled emission, blending
with local learning --- constitutes a practical design pattern for AI systems.
This section strips away the Sheldrake framing and examines the engineering
utility.

% -- 5.1 Population-Level Experience Replay ----------------------------------
\subsection{Population-Level Experience Replay}
\label{sec:applications:replay}

The morphic field is a population-level experience replay buffer where
contributions are weighted by success.  Unlike standard experience
replay~\cite{schaul2015}, which stores individual transitions for a single
agent, the morphic field operates at the \emph{policy level}: it accumulates
entire Q-tables, weighted by evaluation performance, and emits a distilled
initialization for new agents.

This is related to but distinct from:
\begin{itemize}[leftmargin=2em]
    \item \textbf{Knowledge distillation}~\cite{hinton2015}: compresses many
          models into one.  The morphic field does not compress --- it
          accumulates and averages.
    \item \textbf{Federated learning}~\cite{mcmahan2017}: aggregates local
          updates into a global model.  The morphic field is structurally
          similar, but asymmetric --- performance weighting means not all
          agents contribute equally.
    \item \textbf{Population-based training}~\cite{jaderberg2017}: evolves
          hyperparameters across a population.  The morphic field evolves
          policies, not hyperparameters, and includes non-elite contributors
          (at reduced weight).
\end{itemize}

% -- 5.2 Cold-Start Mitigation -----------------------------------------------
\subsection{Cold-Start Mitigation}
\label{sec:applications:coldstart}

The displaced control result --- fresh agents immediately performing 26\%
better --- is a direct solution to the \emph{cold-start problem}.  A new agent
spins up with no history but benefits from a collective memory accumulated by
all prior agents.

Applications include:
\begin{itemize}[leftmargin=2em]
    \item \textbf{Scaling agent swarms}: new workers bootstrap from the
          swarm's accumulated knowledge rather than learning from scratch.
    \item \textbf{Recovery after failure}: a crashed agent's replacement
          inherits the field immediately.
    \item \textbf{Cross-task transfer}: the field serves as a domain prior
          (rather than task-specific weights) that gives new agents a head
          start on related problems.
\end{itemize}

% -- 5.3 Stigmergic Coordination ---------------------------------------------
\subsection{Stigmergic Coordination}
\label{sec:applications:stigmergy}

The morphic field is a form of \emph{stigmergy}~\cite{grasse1959} --- agents
coordinate through a shared environment rather than direct communication.
This is the same principle by which ant colonies optimize
paths~\cite{dorigo1996}, Wikipedia edits accumulate knowledge, and git
repositories coordinate distributed development.

Stigmergic architectures scale better than peer-to-peer messaging because:
(A)~agents do not need to know about each other; (B)~the shared medium
naturally aggregates and attenuates information; (C)~new agents benefit
immediately without explicit onboarding.

\paragraph{Connection to GraphPalace.}
The GraphPalace stigmergic memory engine~\cite{graphpalace2026} uses five-type
pheromone trails for memory retrieval: success, traversal, recency,
exploitation, and exploration pheromones accumulate on graph edges to create
living retrieval landscapes.  The morphic field implements the same principle
in Q-table space: agents deposit performance-weighted Q-values and withdraw
scaled priors.  A natural extension would apply morphic field architecture to
GraphPalace's embedding space, creating a ``learning palace'' where the memory
system itself improves from the agents that use it.

% -- 5.4 Design Pattern Summary ----------------------------------------------
\subsection{Summary: The Morphic Field as a Design Pattern}
\label{sec:applications:summary}

Table~\ref{tab:pattern} summarizes the morphic field as a reusable
architecture.

\begin{table}[ht]
\centering
\caption{The morphic field as a design pattern for multi-agent systems.}
\label{tab:pattern}
\begin{tabular}{@{}ll@{}}
\toprule
\textbf{Component} & \textbf{Description} \\
\midrule
Shared medium    & Accumulated representation (Q-table, embedding, weight
                   matrix) \\
Absorb           & Performance-weighted deposit after each agent's evaluation \\
Emit             & Scaled prior for new agents (strength parameter controls
                   influence) \\
Blend            & Combination with local/genetic initialization (tunable
                   ratio) \\
Performance gate & Failed agents contribute near-zero; solvers contribute
                   proportionally \\
No identity      & The field does not track individual contributors --- only
                   the aggregate \\
\bottomrule
\end{tabular}
\end{table}

The pattern is agnostic to the representation space.  The simulation uses
flat Q-tables; the same absorb/emit/blend cycle could operate on:
\begin{itemize}[leftmargin=2em]
    \item Neural network weight matrices (deep RL agents)
    \item Embedding vectors (semantic memory systems like GraphPalace)
    \item Tool-use strategy distributions (LLM agent systems)
    \item Hierarchical fields (different fields for different skill levels,
          mapping to GraphPalace's wing/room/drawer hierarchy)
\end{itemize}


% =============================================================================
\section{Discussion}
\label{sec:discussion}
% =============================================================================

\subsection{Limitations}
\label{sec:discussion:limitations}

\begin{enumerate}[leftmargin=2em]
    \item \textbf{Tabular Q-learning.}  The simulation uses flat Q-tables on
          a small grid.  Real-world applications would require function
          approximation (deep RL), where the absorb/emit dynamics are more
          complex and potentially unstable.

    \item \textbf{Single maze.}  All experiments use a single procedurally
          generated maze (seed=42).  Generalization across mazes is not
          tested.  The morphic field may overfit to a specific topology.

    \item \textbf{Morphic-only noise.}  The morphic-only condition is noisy
          (39\% improvement vs.\ 86\% for genetic).  Performance-weighted
          averaging is a blunt aggregation mechanism compared to elite
          selection with crossover.

    \item \textbf{No real biology.}  The simulation models artificial agents,
          not rats.  The mapping from Q-learning to animal learning is
          loose at best.

    \item \textbf{Sheldrake's mechanism is unspecified.}  Our formalization
          (shared Q-table) is one of many possible implementations.
          Sheldrake does not specify the mathematical structure of morphic
          fields, so our results are specific to this formalization.

    \item \textbf{Single seed.}  Results are from a single random seed.  A
          full study would require multiple seeds with confidence intervals.
\end{enumerate}

\subsection{Criticisms of Morphic Resonance}
\label{sec:discussion:criticisms}

For completeness, we note the major scientific criticisms of Sheldrake's
hypothesis, which our simulation does not address:

\begin{enumerate}[leftmargin=2em]
    \item \textbf{No known mechanism}: the hypothesis posits influence across
          time and space without an energetic carrier or physical
          mechanism~\cite{rose2006}.
    \item \textbf{Violates established physics}: implies non-local causation
          without known fields, contradicting the standard model.
    \item \textbf{Unfalsifiable boundaries}: critics argue the hypothesis
          can always be adjusted to accommodate negative
          results~\cite{maddox1981}.
    \item \textbf{Alternative explanations}: genetic drift, methodological
          artifacts, and publication bias can explain the cited experimental
          results without invoking new physics.
\end{enumerate}

\subsection{Future Work}
\label{sec:discussion:future}

\begin{enumerate}[leftmargin=2em]
    \item \textbf{Deep RL extension}: replace tabular Q-learning with neural
          network agents and test whether the morphic field (as shared
          weight initialization) scales to high-dimensional state spaces.
    \item \textbf{Multi-maze generalization}: train morphic fields across
          multiple maze topologies to test whether the field captures
          general navigation strategies rather than topology-specific routes.
    \item \textbf{Embedding-space fields}: implement the absorb/emit/blend
          pattern in vector embedding space (e.g., within GraphPalace) to
          create stigmergic memory systems that improve from agent usage.
    \item \textbf{Multi-seed statistical analysis}: run all conditions across
          50+ seeds with bootstrap confidence intervals.
    \item \textbf{Field decay}: implement temporal decay in the morphic field
          (recent agents contribute more) to model Sheldrake's claim that
          recent habits are stronger.
    \item \textbf{Comparison with federated learning}: benchmark the morphic
          field pattern against FedAvg~\cite{mcmahan2017} and other
          aggregation schemes on identical tasks.
\end{enumerate}


% =============================================================================
\section{Conclusion}
\label{sec:conclusion}
% =============================================================================

We have presented a computational model of Sheldrake's morphic resonance
hypothesis, implemented as a multi-agent Q-learning simulation of the
McDougall rat maze experiment.  The simulation demonstrates that:

\begin{enumerate}[leftmargin=2em]
    \item A morphic field, formalized as a performance-weighted shared Q-table,
          transmits useful learning across agent generations with no genetic
          channel (39\% improvement over 60~generations in the morphic-only
          condition).
    \item The displaced control test produces a distinctive signature: fresh
          agents with zero ancestry immediately outperform the control
          baseline by 26\%, benefiting purely from the accumulated collective
          habit.
    \item The hypothesis is formally coherent, the predicted data pattern is
          distinctive and detectable, and the displaced control constitutes a
          viable experimental design.
    \item The morphic field architecture --- absorb, emit, blend --- is a
          practical design pattern for multi-agent knowledge transfer,
          cold-start mitigation, and stigmergic coordination, independent of
          Sheldrake's metaphysics.
\end{enumerate}

The simulation does not prove that morphic resonance exists in nature.  It is
a computational thought experiment that makes the hypothesis precise enough to
test.  The critical next step is empirical: run a McDougall-style experiment
with rigorous displaced controls and look for the signature this simulation
predicts.

Meanwhile, the engineering contribution stands on its own.  The insight that
learning should not die with the agent --- that collective memory can be
infrastructure --- is alive in AI research under names like experience sharing,
population-based training, and federated learning.  The morphic field pattern
provides a clean, minimal implementation of this principle, with a direct path
from biological inspiration to practical multi-agent architecture.

\vspace{1em}
\noindent\textbf{Code availability.}
The simulation is available at \texttt{morphic\_sim.py} in the BUTTERS project
repository.  Dependencies: Python~3, NumPy, Matplotlib.

% =============================================================================
% Bibliography
% =============================================================================

\begin{thebibliography}{99}

\bibitem{sheldrake1981}
R.~Sheldrake, \emph{A New Science of Life: The Hypothesis of Formative
Causation}. London: Blond and Briggs, 1981.

\bibitem{sheldrake1988}
R.~Sheldrake, \emph{The Presence of the Past: Morphic Resonance and the
Habits of Nature}. London: Collins, 1988.

\bibitem{sheldrake2012}
R.~Sheldrake, \emph{The Science Delusion}. London: Coronet, 2012.
(US title: \emph{Science Set Free}.)

\bibitem{maddox1981}
J.~Maddox, ``A book for burning?'' \emph{Nature}, vol.~293, pp.~245--246,
1981.

\bibitem{mcdougall1938}
W.~McDougall, ``Fourth report on a Lamarckian experiment,'' \emph{British
Journal of Psychology}, vol.~28, pp.~321--345, 1938.

\bibitem{agar1954}
W.~E.~Agar, F.~H.~Drummond, O.~W.~Tiegs, and M.~M.~Gunson, ``Fourth (final)
report on a test of McDougall's Lamarckian experiment on the training of
rats,'' \emph{Journal of Experimental Biology}, vol.~31, pp.~307--321, 1954.

\bibitem{koestler1967}
A.~Koestler, \emph{The Ghost in the Machine}. London: Hutchinson, 1967.

\bibitem{colwell2000}
J.~Colwell, S.~Schr\"{o}der, and D.~Sladen, ``The ability to detect unseen
staring: A literature review and empirical tests,'' \emph{British Journal of
Psychology}, vol.~91, pp.~71--85, 2000.

\bibitem{whitehead1929}
A.~N.~Whitehead, \emph{Process and Reality}. New York: Macmillan, 1929.

\bibitem{bohm1980}
D.~Bohm, \emph{Wholeness and the Implicate Order}. London: Routledge, 1980.

\bibitem{jung1959}
C.~G.~Jung, \emph{The Archetypes and the Collective Unconscious}, Collected
Works, vol.~9, part~1. Princeton: Princeton University Press, 1959.

\bibitem{schaul2015}
T.~Schaul, J.~Quan, I.~Antonoglou, and D.~Silver, ``Prioritized experience
replay,'' in \emph{Proc.\ ICLR}, 2016.

\bibitem{hinton2015}
G.~Hinton, O.~Vinyals, and J.~Dean, ``Distilling the knowledge in a neural
network,'' in \emph{NIPS Deep Learning Workshop}, 2015.

\bibitem{jaderberg2017}
M.~Jaderberg, V.~Dalibard, S.~Osindero, W.~K.~Czarnecki, J.~Donahue,
A.~Razavi, O.~Vinyals, T.~Green, I.~Dunning, K.~Simonyan, C.~Fernando,
and K.~Kavukcuoglu, ``Population based training of neural networks,''
\emph{arXiv:1711.09846}, 2017.

\bibitem{mcmahan2017}
H.~B.~McMahan, E.~Moore, D.~Ramage, S.~Hampson, and B.~A.~y~Arcas,
``Communication-efficient learning of deep networks from decentralized data,''
in \emph{Proc.\ AISTATS}, 2017.

\bibitem{grasse1959}
P.-P.~Grass\'{e}, ``La reconstruction du nid et les coordinations
interindividuelles chez \emph{Bellicositermes natalensis} et
\emph{Cubitermes} sp.,'' \emph{Insectes Sociaux}, vol.~6, pp.~41--80, 1959.

\bibitem{dorigo1996}
M.~Dorigo, V.~Maniezzo, and A.~Colorni, ``Ant system: Optimization by a
colony of cooperating agents,'' \emph{IEEE Trans.\ Systems, Man, and
Cybernetics---Part B}, vol.~26, no.~1, pp.~29--41, 1996.

\bibitem{graphpalace2026}
web3guru888, ``GraphPalace: A stigmergic memory palace engine for AI agents,''
Technical Report, April 2026.
\url{https://github.com/web3guru888/GraphPalace}

\bibitem{rose2006}
S.~Rose, \emph{The 21st-Century Brain: Explaining, Mending and Manipulating
the Mind}. London: Vintage, 2006.

\end{thebibliography}

\end{document}
